% Options for packages loaded elsewhere
\PassOptionsToPackage{unicode}{hyperref}
\PassOptionsToPackage{hyphens}{url}
\documentclass[
  ignorenonframetext,
  aspectratio=169,
]{beamer}
\newif\ifbibliography
\usepackage{pgfpages}
\setbeamertemplate{caption}[numbered]
\setbeamertemplate{caption label separator}{: }
\setbeamercolor{caption name}{fg=normal text.fg}
\beamertemplatenavigationsymbolsempty
% remove section numbering
\setbeamertemplate{part page}{
  \centering
  \begin{beamercolorbox}[sep=16pt,center]{part title}
    \usebeamerfont{part title}\insertpart\par
  \end{beamercolorbox}
}
\setbeamertemplate{section page}{
  \centering
  \begin{beamercolorbox}[sep=12pt,center]{section title}
    \usebeamerfont{section title}\insertsection\par
  \end{beamercolorbox}
}
\setbeamertemplate{subsection page}{
  \centering
  \begin{beamercolorbox}[sep=8pt,center]{subsection title}
    \usebeamerfont{subsection title}\insertsubsection\par
  \end{beamercolorbox}
}
% Prevent slide breaks in the middle of a paragraph
\widowpenalties 1 10000
\raggedbottom
\AtBeginPart{
  \frame{\partpage}
}
\AtBeginSection{
  \ifbibliography
  \else
    \frame{\sectionpage}
  \fi
}
\AtBeginSubsection{
  \frame{\subsectionpage}
}
\usepackage{iftex}
\ifPDFTeX
  \usepackage[T1]{fontenc}
  \usepackage[utf8]{inputenc}
  \usepackage{textcomp} % provide euro and other symbols
\else % if luatex or xetex
  \usepackage{unicode-math} % this also loads fontspec
  \defaultfontfeatures{Scale=MatchLowercase}
  \defaultfontfeatures[\rmfamily]{Ligatures=TeX,Scale=1}
\fi
\usepackage{lmodern}
\ifPDFTeX\else
  % xetex/luatex font selection
\fi
% Use upquote if available, for straight quotes in verbatim environments
\IfFileExists{upquote.sty}{\usepackage{upquote}}{}
\IfFileExists{microtype.sty}{% use microtype if available
  \usepackage[]{microtype}
  \UseMicrotypeSet[protrusion]{basicmath} % disable protrusion for tt fonts
}{}
\makeatletter
\@ifundefined{KOMAClassName}{% if non-KOMA class
  \IfFileExists{parskip.sty}{%
    \usepackage{parskip}
  }{% else
    \setlength{\parindent}{0pt}
    \setlength{\parskip}{6pt plus 2pt minus 1pt}}
}{% if KOMA class
  \KOMAoptions{parskip=half}}
\makeatother
\usepackage{longtable,booktabs,array}
\usepackage{caption}
\captionsetup[table]{skip=6pt}
\newcounter{none} % for unnumbered tables
\usepackage{calc} % for calculating minipage widths
\usepackage{caption}
% Make caption package work with longtable
\makeatletter
\def\fnum@table{\tablename~\thetable}
\makeatother
\setlength{\emergencystretch}{3em} % prevent overfull lines
\providecommand{\tightlist}{%
  \setlength{\itemsep}{0pt}\setlength{\parskip}{0pt}}
% Auto-generated by infrastructure/rendering/_pdf_latex_helpers.py::extract_math_font_preamble.
% Minimal header passed to Pandoc via -H so Beamer slides pick up the same math font
% as the combined-PDF path. Adjust the manuscript's preamble.md to change the font globally.
\usepackage{unicode-math}
\setmathfont{latinmodern-math.otf}

% Manuscript macro/package fallbacks (newcommand -> providecommand).
% Core mathematics
\usepackage{amsmath}
\usepackage{amssymb}
% Document layout
% Tables
\usepackage{booktabs}
% Caption styling (booktabs already loaded above). Small caption font with a
% bold label ("Figure 1", "Table 2") gives the math/figure-heavy layout a
% consistent, journal-like caption rhythm without fighting pandoc-crossref —
% pandoc emits standard \caption{}, and caption/captionsetup only restyle it.
% ── Theorem environments ─────────────────────────────────────────────────
% amsthm is loaded above. The FedGVI<->Friston bridge is stated as a small
% number of theorems/definitions/lemmas/propositions/corollaries; sharing one
% counter (definition/lemma/proposition/corollary all step the `theorem`
% counter) gives a single monotone numbering 1, 2, 3, ... across all kinds, so
% authors can reference them as Theorem 1, Definition 2, Corollary 3 without a
% per-kind counter drifting out of order. Plain style for theorem-like results,
% definition style (upright body) for definitions.
% (algorithm2e intentionally not loaded: this manuscript states results as
% numbered amsthm Theorems/Definitions, not algorithm2e pseudocode.)
% Code listings
\usepackage{listings}
% Typography and formatting
% (siunitx intentionally not loaded: no \SI/\num units appear in this manuscript.)
% Cross-references and citations
% (cleveref and titlesec intentionally not loaded: unavailable in this TeX tree
% and not required. Cross-references resolve through pandoc-crossref's own
% \ref-style names — no \cref appears — and section heading styling falls back
% to the pandoc/LaTeX defaults.)
% ── TeX Live Basic-compatible mono font for code listings ────────────
% Latin Modern Mono is bundled with TeX Live Basic and keeps the manuscript
% renderable on minimal TeX installations. Mathematical Unicode coverage is
% provided separately by Latin Modern Math below, so the code font does not
% need a private JuliaMono installation.
%
% Use the TeX font filename rather than the system-family name: the minimal
% TeX Live fontconfig database may not register the OpenType family even though
% kpsewhich can resolve the bundled files.
% Math font for unicode-math: Latin Modern Math (TeX Live) has full BMP
% coverage including U+2223 (\mid), U+226A/226B (\ll/\gg), and the Greek/
% blackboard letters used in equations. Without an explicit \setmathfont,
% unicode-math falls back to lmroman text font which lacks several glyphs
% and emits "Missing character" warnings on every \mid in math mode.

% Slide layout defaults for warning-clean scientific decks.
\usepackage{etoolbox}
\IfFileExists{xurl.sty}{\usepackage{xurl}}{}
\IfFileExists{seqsplit.sty}{\usepackage{seqsplit}}{\newcommand{\seqsplit}[1]{#1}}
\protected\def\breakseq#1{\seqsplit{#1}}
\protected\def\breaktt#1{\begingroup\ttfamily\seqsplit{#1}\endgroup}
\setlength{\emergencystretch}{6em}
\tolerance=5000
\hbadness=10000
\hfuzz=1pt
\setlength{\tabcolsep}{2pt}
\AtBeginEnvironment{longtable}{\tiny\renewcommand{\arraystretch}{0.86}\setlength{\tabcolsep}{1pt}}
\AtBeginEnvironment{tabular}{\tiny\renewcommand{\arraystretch}{0.86}\setlength{\tabcolsep}{1pt}}
\AtBeginEnvironment{equation}{\tiny}
\AtBeginEnvironment{equation*}{\tiny}
\AtBeginEnvironment{align}{\tiny}
\AtBeginEnvironment{align*}{\tiny}
\AtBeginEnvironment{itemize}{\footnotesize}
\AtBeginEnvironment{enumerate}{\footnotesize}
\AtBeginEnvironment{description}{\footnotesize}
\setbeamerfont{caption}{size=\tiny}
\setbeamerfont{caption name}{size=\tiny}
\setbeamerfont{normal text}{size=\small}
\setbeamerfont{frametitle}{size=\small}
\setbeamerfont{section title}{size=\footnotesize}
\setbeamerfont{subsection title}{size=\footnotesize}
\setbeamertemplate{section page}{%
  \centering
  \begin{beamercolorbox}[sep=12pt,center,wd=\paperwidth]{section title}
    \parbox{0.86\paperwidth}{\centering\usebeamerfont{section title}\insertsection\par}
  \end{beamercolorbox}
}
\setbeamertemplate{subsection page}{%
  \centering
  \begin{beamercolorbox}[sep=8pt,center,wd=\paperwidth]{subsection title}
    \parbox{0.86\paperwidth}{\centering\usebeamerfont{subsection title}\insertsubsection\par}
  \end{beamercolorbox}
}
\setlength{\abovecaptionskip}{2pt}
\setlength{\belowcaptionskip}{0pt}

% Natbib and cross-reference fallbacks — slides don't load natbib
% or cleveref, but combined-PDF manuscript prose may emit these
% commands. The fallback renders citations as a bracketed key list
% and cross-references as detokenized labels so slides don't fail on
% undefined control sequences or raw underscores. \providecommand is
% a no-op if packages load later.
\providecommand{\citep}[1]{[#1]}
\providecommand{\citet}[1]{#1}
\providecommand{\citealp}[1]{#1}
\providecommand{\citeauthor}[1]{#1}
\providecommand{\citeyear}[1]{#1}
\providecommand{\cref}[1]{\texttt{\detokenize{#1}}}
\providecommand{\Cref}[1]{\texttt{\detokenize{#1}}}
\providecommand{\crefrange}[2]{\texttt{\detokenize{#1}}--\texttt{\detokenize{#2}}}
\providecommand{\Crefrange}[2]{\texttt{\detokenize{#1}}--\texttt{\detokenize{#2}}}
\providecommand{\paragraph}[1]{\textbf{#1}\ }

% Opt-in accessible presentation profile.
\setbeamerfont{normal text}{size*={20pt}{24pt}}
\setbeamerfont{frametitle}{size*={28pt}{32pt}}
\setbeamerfont{section title}{size*={28pt}{32pt}}
\setbeamerfont{subsection title}{size*={28pt}{32pt}}
\setbeamerfont{caption}{size*={16pt}{19pt}}
\setbeamerfont{caption name}{size*={16pt}{19pt}}
\setbeamerfont{itemize/enumerate body}{size*={20pt}{24pt}}
\setbeamerfont{itemize/enumerate subbody}{size*={20pt}{24pt}}
\setbeamerfont{itemize/enumerate subsubbody}{size*={20pt}{24pt}}
\setbeamerfont{itemize item}{size*={20pt}{24pt}}
\setbeamerfont{itemize subitem}{size*={20pt}{24pt}}
\setbeamerfont{itemize subsubitem}{size*={20pt}{24pt}}
\setbeamerfont{description body}{size*={20pt}{24pt}}
\setbeamerfont{description item}{size*={20pt}{24pt}}
\setbeamerfont{quote}{size*={20pt}{24pt}}
\setbeamerfont{quotation}{size*={20pt}{24pt}}
\AtBeginEnvironment{longtable}{\fontsize{20pt}{24pt}\selectfont}
\AtBeginEnvironment{tabular}{\fontsize{20pt}{24pt}\selectfont}
\AtBeginEnvironment{equation}{\fontsize{20pt}{24pt}\selectfont}
\AtBeginEnvironment{equation*}{\fontsize{20pt}{24pt}\selectfont}
\AtBeginEnvironment{align}{\fontsize{20pt}{24pt}\selectfont}
\AtBeginEnvironment{align*}{\fontsize{20pt}{24pt}\selectfont}
\AtBeginEnvironment{itemize}{\fontsize{20pt}{24pt}\selectfont}
\AtBeginEnvironment{enumerate}{\fontsize{20pt}{24pt}\selectfont}
\AtBeginEnvironment{description}{\fontsize{20pt}{24pt}\selectfont}
\AtBeginDocument{\usebeamerfont{normal text}}
\newlength{\AFSlideTableWidth}
% The fixed footer reservation is calibrated with the semantic seven-line regular-body budget.
\setbeamertemplate{footline}{%
  \leavevmode\vbox{%
    \hbox{%
      \begin{beamercolorbox}[wd=\paperwidth,ht=16pt,dp=3pt,center]{author in head/foot}%
        \fontsize{16pt}{19pt}\selectfont Untagged PDF derivative \textbar\ \href{../web/index.html}{HTML reader}%
      \end{beamercolorbox}%
    }%
    \vskip6pt%
  }%
}
% Accessible footer reserved height: 25pt.

\newtheorem{proposition}{Proposition}
\newtheorem{hypothesis}{Hypothesis}
\newtheorem{remark}[theorem]{Remark}
\newtheorem{axiom}{Axiom}
\newtheorem{property}{Property}
\usepackage{bookmark}
\IfFileExists{xurl.sty}{\usepackage{xurl}}{} % add URL line breaks if available
\urlstyle{same}
% fallback for those not using the hyperref driver hyperxmp:
\makeatletter
\@ifundefined{xmpquote}{\newcommand{\xmpquote}[1]{#1}}{}
\makeatother
\hypersetup{
  hidelinks,
  pdfcreator={LaTeX via pandoc}}

\author{\texorpdfstring{}{}}
\date{}

\begin{document}

\begin{frame}{Limitations and claim boundaries}
\protect\phantomsection\label{sec:limitations}
The boundaries here define the contribution. The central goal is to show
--- concretely and testably --- that the categorical
generalized-variational construction inspired by FedGVI (Mildner et al.
2025) has standard-Bayes client limits and a project-local
log-linear-pool server identity.
\end{frame}

\begin{frame}{Limitations and claim\ldots{} (part 2)}
Under the explicit bridge in sec.~5.8, the
latter specializes Friston et al.'s Eq. 7 message-combination term
(Friston et al. 2024), not the complete source protocol; behavior away
from those limits is evaluated only under declared simulation
conditions.
\end{frame}

\begin{frame}{Limitations and claim\ldots{} (part 3)}
The exact identities are formal; performance and deployment claims
remain conditional. Items beyond that boundary are named as future work.
\end{frame}

\begin{frame}{Three robustness axes: theorem, heuristic, and objective}
\protect\phantomsection\label{sec:limitations-axes}
The authoritative definitions and claim owners are in
sec.~3.3; this section records only what still
cannot transfer. A source-conditional client theorem does not certify
either server rule,
\end{frame}

\begin{frame}{Three robustness axes (part 2)}
a server recovery identity does not supply an objective or
estimator-level influence bound, and objective descent plus raw-weight
control does not imply truth recovery or peak-accuracy dominance.
Likewise,
\end{frame}

\begin{frame}{Three robustness axes (part 3)}
conditional accuracy contrasts cannot confer a theorem on the heuristic.
The evidence map in fig.~6 makes that
no-transfer boundary and its nested replication units explicit.
\end{frame}

\begin{frame}{Three robustness axes (part 4)}
The remaining methodological problem is a server objective that combines
competitive conditional performance with an appropriate robustness
guarantee. Related robust-federation work motivates the target
(Blanchard et al. 2017; Li et al. 2022) but does not certify either
project server rule.
\end{frame}

\begin{frame}{Three robustness axes (part 5)}
Any such result must also survive beyond the present categorical model,
declared attack grid, and single-host execution setting; future work and
extended methods remain separate (sec.~8.17,
sec.~13).
\end{frame}

\begin{frame}{Scope boundaries that the evidence does not cross}
\protect\phantomsection\label{sec:limitations-scope}
\textbf{The bridge is a recasting, not an upstream claim.} Friston et
al. (Friston et al. 2024) supply a belief-sharing operator for agents
with a shared world model; Mildner et al.
\end{frame}

\begin{frame}{Scope boundaries that the\ldots{} (part 2)}
(Mildner et al. 2025) supply robust federated generalized variational
inference. Neither paper claims the other.
\end{frame}

\begin{frame}{Scope boundaries that the\ldots{} (part 3)}
The contribution here is to recast the former as the tested KL/NLL
zero-robustness recovery corner of the latter-inspired construction,
then measure changes under robust losses and server reweighting.
\end{frame}

\begin{frame}{Scope boundaries that the\ldots{} (part 4)}
\textbf{Historical sources are conceptual, not formal support.} The
pre-modern sources added in sec.~8.9 support a
genealogy of expectation, inverse inference, utility, and collective
judgment (Huygens 1657; Bayes 1763; Laplace 1774; Bernoulli 1738;
Condorcet 1785). They do not support claims about KL, NLL,
\end{frame}

\begin{frame}[fragile]{Scope boundaries that the\ldots{} (part 5)}
product-of-experts training, FedGVI, or \texttt{robust\_aggregate}.

Those claims rest only on the modern cited formalism and the tests
reported in sec.~7.1.
\end{frame}

\begin{frame}{Scope boundaries that the\ldots{} (part 6)}
\textbf{Single-machine federation.} Federation is validated with queue
transport, a single-machine OS-process helper, and loopback TCP: agents
serialize beliefs, the server aggregates, and consensus is broadcast
back without changing the mathematics.
\end{frame}

\begin{frame}{Scope boundaries that the\ldots{} (part 7)}
The socket path exercises frame integrity and file-backed
digest-verified replay validation; an optional SQLite round-ID guard
survives local process restarts.
\end{frame}

\begin{frame}{Scope boundaries that the\ldots{} (part 8)}
These controls do not substitute for cross-host deployment,
identity-bound mTLS, a shared multi-host replay domain, discovery,
long-running worker orchestration, or fault tolerance; those steps are
scoped as future work (sec.~8.17).
\end{frame}

\begin{frame}{Scope boundaries that the\ldots{} (part 9)}
\textbf{Source-scale posterior FedGVI classification remains open.} The
displayed NumPy logistic-regression baseline is an exploratory
point-estimate comparison, not the source paper's posterior-uncertainty
BNN experiment (Mildner et al. 2025).
\end{frame}

\begin{frame}{Scope boundaries that the\ldots{} (part 10)}
Its operating point was selected within the evaluated synthetic sweep to
expose a mid-contamination margin (fig.~18).
\end{frame}

\begin{frame}{Scope boundaries that the\ldots{} (part 11)}
The curves lose reliable ordering at the highest rate. The source
theorem, not this selected finite sweep, carries the source-conditional
robustness claim.
\end{frame}

\begin{frame}{Scope boundaries that the\ldots{} (part 12)}
\textbf{Classification baselines are point estimates, not full
posteriors.} The NumPy logistic-regression baseline and the PyTorch MLP
complement both use point-estimate weights; no posterior covariance is
computed.
\end{frame}

\begin{frame}[fragile]{Scope boundaries that the\ldots{} (part 13)}
A genuine mean-field variational family over the
weights---diagonal-Gaussian \(q(w)\) with a closed-form KL and a
Monte-Carlo ELBO---is implemented by the tested \texttt{VariationalMLP}
primitive in the \texttt{bnn\_variational\_torch} module. It recovers
the point-estimate net exactly as its posterior variance vanishes.
\end{frame}

\begin{frame}{Scope boundaries that the\ldots{} (part 14)}
The source-comparable training and evaluation lane remains unexecuted:
leakage-free calibration, source-dataset parity, posterior-family
comparisons, and source-scale compute are still open.
\end{frame}

\begin{frame}{Scope boundaries that the\ldots{} (part 15)}
\textbf{Hierarchical depth does not by itself improve the base task.}
The two- and three-level stacks (sec.~14.9,
sec.~14.11) run alternating L1/L2(/L3) inference
end-to-end and resolve their added context latents above chance. On the
shared location task they do not beat the flat baseline:
\end{frame}

\begin{frame}{Scope boundaries that the\ldots{} (part 16)}
the paired location-accuracy gap is a small, statistically reliable
negative gap at the reported seed count
(fig.~28).
\end{frame}

\begin{frame}{Scope boundaries that the\ldots{} (part 17)}
Depth is therefore validated as executable and consensus-preserving, not
as an accuracy improvement; whether a richer policy-and-horizon task
family rewards hierarchy is future work
(sec.~8.20).
\end{frame}

\begin{frame}{Scope boundaries that the\ldots{} (part 18)}
\textbf{Discrete POMDP only---continuous or hybrid state spaces are not
addressed.} This work is discrete-categorical only, matching the
community's worked POMDP example (Da Costa et al. 2020; Friston et al.
2024).
\end{frame}

\begin{frame}{Scope boundaries that the\ldots{} (part 19)}
The limit-as-proof contract is validated off the categorical case only
for a one-dimensional Gaussian-mean conjugate slice; continuous-state
active inference and Gaussian belief-sharing colonies remain untested.
\end{frame}

\begin{frame}{Scope boundaries that the\ldots{} (part 20)}
\textbf{Temperature and divergence calibration are fixed, not learned.}
Coarsened posterior and safe-Bayes theory make clear that the
learning-rate/temperature is part of the inference rule under
misspecification (Miller and Dunson 2018; Grünwald 2012; Kleijn and
Vaart 2012).
\end{frame}

\begin{frame}{Scope boundaries that the\ldots{} (part 21)}
This manuscript validates the configured losses and divergences, but it
does not learn an optimal coarsening radius, temperature, or divergence
schedule across agents.
\end{frame}

\begin{frame}{Scope boundaries that the\ldots{} (part 22)}
\textbf{Real multi-machine federation, networking, and privacy
cryptography.} Federation here is \emph{mathematical} (factor
aggregation), not infrastructural--- unlike communication-efficient or
Byzantine-robust federated learning (McMahan et al. 2017; Blanchard et
al. 2017).
\end{frame}

\begin{frame}{Scope boundaries that the\ldots{} (part 23)}
\textbf{New linguistic theory.} The language-acquisition study
(sec.~7.3) reproduces the Dirichlet count
mechanism (eq.~18) mechanically; it does not
extend the linguistics.
\end{frame}

\begin{frame}{What the statistics can and cannot claim}
\protect\phantomsection\label{sec:limitations-stats}
The verdict is a paired comparison at a single high contamination rate
(sec.~7.5.2, tbl.~7). The
design concentrates power where the effect is largest.
\end{frame}

\begin{frame}{What the statistics can and\ldots{} (part 2)}
It therefore identifies \emph{where} robustness pays off, but it does
not characterize contamination as a continuous function.
\end{frame}

\begin{frame}{What the statistics can and\ldots{} (part 3)}
The per-rate effect and inference projections
(tbls.~5, 6)
describe the remaining sweep.
\end{frame}

\begin{frame}{What the statistics can and\ldots{} (part 4)}
Their inferential claims do not exceed the declared BH-FDR control
(Benjamini and Hochberg 1995).
\end{frame}

\begin{frame}{What the statistics can and\ldots{} (part 5)}
BH-FDR controls expected false discovery proportion within the declared
family, not the chance of any false positive. The matched-pairs Wilcoxon
tests are rank tests under paired-difference assumptions (Wilcoxon 1945;
Fay and Proschan 2010), not assumption-free proofs about raw means.
\end{frame}

\begin{frame}{What the statistics can and\ldots{} (part 6)}
Confidence intervals are percentile bootstrap (Efron and Tibshirani
1993), not analytic, and inherit that method's small-sample caveats at
the lowest trial counts.
\end{frame}

\begin{frame}{What the statistics can and\ldots{} (part 7)}
The power values are observed-effect design approximations useful for
confirmatory sample-size planning; they are not independent evidence for
the verdict and do not cover model-specification uncertainty outside the
seeded simulation harness (Wasserstein and Lazar 2016).
\end{frame}

\end{document}
